% GENERATED FILE. Edit canonical lessons, then run npm run generate:books.
% canonical-source-hash: fnv1a64:c6ebf9a7d7af98bf
% canonical-lessons: BN-C362-monojog, BN-C362-bibhranto, BN-C362-upolobdhi-kora, BN-C362-bodh-kora, BN-C362-sposhto

\chapter{Making Sense of It}
\label{ch:bn-making-sense-of-it}
\hlchaptermodality{362}

\begin{chapteropening}
\textbf{By the end of this chapter:} \emph{I can say when something is clear to me and when I am lost.}

Complete the last lesson of chapter 362: I can say when something is clear to me and when I am lost.
\end{chapteropening}

\section[\texorpdfstring{\bn{মনোযোগ} (mônojog)}{মনোযোগ (mônojog)}]{\bn{মনোযোগ} (mônojog) — attention; concentration}

\label{lesson:BN-C362-monojog}



\begin{quote}
Before the new one: say the Bengali for a university degree, then the Bengali for to tutor.
\end{quote}

\subsection*{You'll want to know: \bn{মনোযোগ}}
\textbf{\bn{মনোযোগ}} — \emph{mônojog} — \textquotedblleft{}attention; concentration\textquotedblright{}.

\begin{quote}
\bn{মনোযোগ} \bn{দিয়ে} \bn{শোনো।}
\end{quote}

\begin{quote}
\emph{mônojog diye shono.}
\end{quote}

\begin{quote}
Listen carefully.
\end{quote}

\begin{cousinweb}
\textbf{\bn{মনোযোগ}} (\emph{mônojog}) is a learned compound of Sanskrit \emph{manas}, mind, and \emph{yoga}, joining (Monier-Williams): joining your mind to the task. The compound itself is not in the dictionary. \bn{মনোযোগ} \bn{দিয়ে} means 'carefully, with attention'.
\end{cousinweb}

\begin{grammarlens}[title={What you've built}]
One more word for giving reasons and opinions.
\end{grammarlens}

\subsection*{Your turn}
\begin{itemize}
\raggedright
  \item \emph{Say it:} \emph{mônojog}
  \item \emph{Say it:} \emph{mônojog}, once more
  \item \emph{Say it:} \emph{mônojog diye shono.} — Listen carefully.
  \item \emph{Recall:} say the Bengali for a university degree, then the Bengali for to tutor, then say \emph{mônojog} again
\end{itemize}

\begin{tcolorbox}[breakable,colback=teal!4,colframe=teal!35!black,title={Before you move on}]
What does \bn{মনোযোগ} mean? (\textquotedblleft{}Attention; concentration\textquotedblright{}.) Say it once more.
\end{tcolorbox}
\section[\texorpdfstring{\bn{বিভ্রান্ত} (bibhrāntô)}{বিভ্রান্ত (bibhrāntô)}]{\bn{বিভ্রান্ত} (bibhrāntô) — confused; misled}

\label{lesson:BN-C362-bibhranto}



\begin{quote}
Before the new one: say the Bengali for to tutor, then the Bengali for concentration.
\end{quote}

\subsection*{You'll want to know: \bn{বিভ্রান্ত}}
\textbf{\bn{বিভ্রান্ত}} — \emph{bibhrāntô} — \textquotedblleft{}confused; misled\textquotedblright{}.

\begin{quote}
\bn{ভুল} \bn{খবর} \bn{শুনে} \bn{অনেকে} \bn{বিভ্রান্ত} \bn{হল।}
\end{quote}

\begin{quote}
\emph{bhul khôbôr shune ôneke bibhrāntô hôlô.}
\end{quote}

\begin{quote}
Many people were misled by the false news.
\end{quote}

\begin{cousinweb}
\textbf{\bn{বিভ্রান্ত}} (\emph{bibhrāntô}) is Sanskrit \emph{vibhrānta}, wandering about, and also confused or bewildered (Monier-Williams): a mind that has lost its way.
\end{cousinweb}

\begin{grammarlens}[title={What you've built}]
One more word for giving reasons and opinions.
\end{grammarlens}

\subsection*{Your turn}
\begin{itemize}
\raggedright
  \item \emph{Say it:} \emph{bibhrāntô}
  \item \emph{Say it:} \emph{bibhrāntô}, once more
  \item \emph{Say it:} \emph{bhul khôbôr shune ôneke bibhrāntô hôlô.} — Many people were misled by the false news.
  \item \emph{Recall:} say the Bengali for to tutor, then the Bengali for concentration, then say \emph{bibhrāntô} again
\end{itemize}

\begin{tcolorbox}[breakable,colback=teal!4,colframe=teal!35!black,title={Before you move on}]
What does \bn{বিভ্রান্ত} mean? (\textquotedblleft{}Confused; misled\textquotedblright{}.) Say it once more.
\end{tcolorbox}
\section[\texorpdfstring{\bn{উপলব্ধি} \bn{করা} (upôlôbdhi kôrā)}{উপলব্ধি করা (upôlôbdhi kôrā)}]{\bn{উপলব্ধি} \bn{করা} (upôlôbdhi kôrā) — to realise; to grasp}

\label{lesson:BN-C362-upolobdhi-kora}



\begin{quote}
Before the new one: say the Bengali for concentration, then the Bengali for misled.
\end{quote}

\subsection*{You'll want to know: \bn{উপলব্ধি} \bn{করা}}
\textbf{\bn{উপলব্ধি} \bn{করা}} — \emph{upôlôbdhi kôrā} — \textquotedblleft{}to realise; to grasp\textquotedblright{}.

\begin{quote}
\bn{পরে} \bn{আমি} \bn{আমার} \bn{ভুল} \bn{উপলব্ধি} \bn{করলাম।}
\end{quote}

\begin{quote}
\emph{pôre āmi āmār bhul upôlôbdhi kôrlām.}
\end{quote}

\begin{quote}
Later I realised my mistake.
\end{quote}

Past (he/she …): \bn{তিনি} \bn{কথাটার} \bn{গুরুত্ব} \bn{উপলব্ধি} \bn{করলেন}, \emph{tini kôthāṭār gurutvô upôlôbdhi kôrlen}.

\begin{cousinweb}
\textbf{\bn{উপলব্ধি} \bn{করা}} (\emph{upôlôbdhi kôrā}) is Sanskrit \emph{upalabdhi}, becoming aware, understanding (Monier-Williams), with \bn{করা}: to come to see a thing clearly.
\end{cousinweb}

\begin{grammarlens}[title={What you've built}]
One more word for giving reasons and opinions.
\end{grammarlens}

\subsection*{Your turn}
\begin{itemize}
\raggedright
  \item \emph{Say it:} \emph{upôlôbdhi kôrā}
  \item \emph{Say it:} \emph{upôlôbdhi kôrā}, once more
  \item \emph{Say it:} \emph{pôre āmi āmār bhul upôlôbdhi kôrlām.} — Later I realised my mistake.
  \item \emph{Recall:} say the Bengali for concentration, then the Bengali for misled, then say \emph{upôlôbdhi kôrā} again
\end{itemize}

\begin{tcolorbox}[breakable,colback=teal!4,colframe=teal!35!black,title={Before you move on}]
What does \bn{উপলব্ধি} \bn{করা} mean? (\textquotedblleft{}To realise; to grasp\textquotedblright{}.) Say it once more.
\end{tcolorbox}
\section[\texorpdfstring{\bn{বোধ} \bn{করা} (bodh kôrā)}{বোধ করা (bodh kôrā)}]{\bn{বোধ} \bn{করা} (bodh kôrā) — to feel (a state); to sense}

\label{lesson:BN-C362-bodh-kora}



\begin{quote}
Before the new one: say the Bengali for misled, then the Bengali for to realise.
\end{quote}

\subsection*{You'll want to know: \bn{বোধ} \bn{করা}}
\textbf{\bn{বোধ} \bn{করা}} — \emph{bodh kôrā} — \textquotedblleft{}to feel (a state); to sense\textquotedblright{}.

\begin{quote}
\bn{সারাদিন} \bn{কাজ} \bn{করে} \bn{আমি} \bn{ক্লান্ত} \bn{বোধ} \bn{করছি।}
\end{quote}

\begin{quote}
\emph{sārādin kāj kôre āmi klāntô bodh kôrchhi.}
\end{quote}

\begin{quote}
After working all day I feel tired.
\end{quote}

Past (he/she …): \bn{তিনি} \bn{ক্লান্ত} \bn{বোধ} \bn{করলেন}, \emph{tini klāntô bodh kôrlen}.

\begin{cousinweb}
\textbf{\bn{বোধ} \bn{করা}} (\emph{bodh kôrā}) is Sanskrit \emph{bodha}, knowing or understanding (Monier-Williams), with \bn{করা}: to become aware of a state in yourself, tiredness or relief. It is a little more formal than \bn{মনে} \bn{হওয়া}.
\end{cousinweb}

\begin{grammarlens}[title={What you've built}]
One more word for giving reasons and opinions.
\end{grammarlens}

\subsection*{Your turn}
\begin{itemize}
\raggedright
  \item \emph{Say it:} \emph{bodh kôrā}
  \item \emph{Say it:} \emph{bodh kôrā}, once more
  \item \emph{Say it:} \emph{sārādin kāj kôre āmi klāntô bodh kôrchhi.} — After working all day I feel tired.
  \item \emph{Recall:} say the Bengali for misled, then the Bengali for to realise, then say \emph{bodh kôrā} again
\end{itemize}

\begin{tcolorbox}[breakable,colback=teal!4,colframe=teal!35!black,title={Before you move on}]
What does \bn{বোধ} \bn{করা} mean? (\textquotedblleft{}To feel (a state); to sense\textquotedblright{}.) Say it once more.
\end{tcolorbox}
\section[\texorpdfstring{\bn{স্পষ্ট} (spôshṭô)}{স্পষ্ট (spôshṭô)}]{\bn{স্পষ্ট} (spôshṭô) — clear; distinct}

\label{lesson:BN-C362-sposhto}



\begin{quote}
Before the new one: say the Bengali for to realise, then the Bengali for to sense.
\end{quote}

\subsection*{You'll want to know: \bn{স্পষ্ট}}
\textbf{\bn{স্পষ্ট}} — \emph{spôshṭô} — \textquotedblleft{}clear; distinct\textquotedblright{}.

\begin{quote}
\bn{তোমার} \bn{কথাটা} \bn{আমার} \bn{কাছে} \bn{স্পষ্ট} \bn{নয়।}
\end{quote}

\begin{quote}
\emph{tomār kôthāṭā āmār kāchhe spôshṭô nôy.}
\end{quote}

\begin{quote}
What you said isn't clear to me.
\end{quote}

\begin{cousinweb}
\textbf{\bn{স্পষ্ট}} (\emph{spôshṭô}) is Sanskrit \emph{spaṣṭa}, clearly seen, distinct (Monier-Williams). Speech, writing or a picture can be \bn{স্পষ্ট}.
\end{cousinweb}

\begin{grammarlens}[title={What you've built}]
One more word for giving reasons and opinions.
\end{grammarlens}

\subsection*{Your turn}
\begin{itemize}
\raggedright
  \item \emph{Say it:} \emph{spôshṭô}
  \item \emph{Say it:} \emph{spôshṭô}, once more
  \item \emph{Say it:} \emph{tomār kôthāṭā āmār kāchhe spôshṭô nôy.} — What you said isn't clear to me.
  \item \emph{Recall:} say the Bengali for to realise, then the Bengali for to sense, then say \emph{spôshṭô} again
\end{itemize}

\begin{tcolorbox}[breakable,colback=teal!4,colframe=teal!35!black,title={Before you move on}]
What does \bn{স্পষ্ট} mean? (\textquotedblleft{}Clear; distinct\textquotedblright{}.) Say it once more.
\end{tcolorbox}
